\begin{definition}
Let $g\in\mathsf{IncSeq}$, $k\in\mathbb N$, $f\in\mathsf{IncSeq}$, and
$l\in\mathbb N$.  Put
\[
G(n)=\mathsf{OrbitPoint}(g,k,n)=g^{[n]}(k).
\]
The total block-map value $\mathsf{BlockSigma}(g,k,f,l)$ is defined as
follows.  If
\[
l<\mathsf{OrbitPoint}(g,k,0),
\]
then it is $l$.  Otherwise, if there exists an $n$ such that
\[
l<\mathsf{OrbitPoint}(g,k,n+1)=G(n+1),
\]
let $n_l$ be the least such $n$ and set
\[
\mathsf{BlockSigma}(g,k,f,l)=g^{[f(n_l)-n_l]}(l).
\]
If no such $n$ exists, set
\[
\mathsf{BlockSigma}(g,k,f,l)=l.
\]

When $k$ is the least moved point of $g$, meaning
$k<g(k)$ and $g(j)=j$ for every $j<k$, the orbit is unbounded and every
$l\geq G(0)$ lies in a unique block $G(n)\leq l<G(n+1)$.  Thus the final
fallback is unreachable, and the definition specializes to the paper's
two-case block map
\[
\sigma_f(l)=
\begin{cases}
l,&l<G(0),\\
g^{[f(n)-n]}(l),&G(n)\leq l<G(n+1)\text{ for the unique }n.
\end{cases}
\]
This is the block map used in \texttt{BlockSigmaIntertwines} and
\texttt{EmapLemma}.
\end{definition}
